% Generated draft: requires mathematical and source review.
\documentclass[11pt]{article}
\usepackage[margin=25mm]{geometry}
\usepackage{fontspec}
\setmainfont{Latin Modern Roman}
\usepackage{amsmath,amssymb,mathtools,unicode-math}
\setmathfont{Latin Modern Math}
\usepackage{enumitem,needspace,xurl,hyperref}
\setlength{\parindent}{0pt}
\setlength{\parskip}{5pt}
\setlength{\emergencystretch}{3em}
\newcommand{\N}{\mathbb{N}}
\newcommand{\Z}{\mathbb{Z}}
\newcommand{\Q}{\mathbb{Q}}
\newcommand{\R}{\mathbb{R}}
\newcommand{\C}{\mathbb{C}}
\newcommand{\F}{\mathbb{F}}
\newcommand{\A}{\mathbb{A}}
\newcommand{\PP}{\mathbb{P}}
\newcommand{\E}{\mathbb{E}}
\newcommand{\eps}{\varepsilon}
\newcommand{\dd}{\,\mathrm{d}}
\begin{document}
\section*{UnsolvedMath ID: 6700065}
\textbf{Source problem code:} catalogue entry
\subsection{Definitions and mathematical statement}
Let \(X^n\) be the source's compact orientable Riemannian manifold and \(f_0:X\to T^n\) a continuous map of nonzero degree. In the usual closed-manifold convention, \((f_0)_*[X]=d[T^n]\) with \(d\in\Z\setminus\{0\}\). Let \(h_i\in H_1(T^n;\Z)\), \(i\ge0\), collectively generate that group, and let \(\varepsilon_i>0\) with \(\varepsilon_i\to0\).

The preceding construction of the regularized spaces \(X_i,X'_i\) and minimizing hypersurfaces \(Y_i\subset X'_i\) is absent. Conditional on that construction and its metric normalizations, the shrinking conjecture asserts
\[
\operatorname{diam}(X'_i)\longrightarrow0,
\qquad
\exists i_0\ \forall i\ge i_0,\quad \operatorname{Sing}(Y_i)=\varnothing.
\]
It further asserts near-isometries to flat \(n\)-tori: there are \(\lambda_i\ge1\), \(\lambda_i\to1\), such that each \(Q\in\{X_i,X'_i\}\) admits a homeomorphism \(\phi:Q\to T_{i,Q}\) to a flat torus with
\[
\lambda_i^{-1}d_Q(x,y)\le d_{T_{i,Q}}(\phi(x),\phi(y))
\le\lambda_i d_Q(x,y)\quad(x,y\in Q).
\]
Separate torus symbols avoid the extraction's reuse of \(Y_i\) for both an \((n-1)\)-dimensional hypersurface and an \(n\)-torus; whether a common torus is required is unresolved.

The source says this mechanism would imply flatness of overtorical manifolds, meaning manifolds with a nonzero-degree map to \(T^n\), under \(\operatorname{Sc}\ge0\). It invokes scalar-curvature semicontinuity under Lipschitz limits, referring to a \(C^0\) argument in [49], and suggests that arguments in [86,115] might prove the mechanism. These are reported implications and suggestions, not a supplied proof. In particular, \(\operatorname{Sc}(X)\ge0\) is present in the asserted consequence but not explicitly in the extracted initial hypothesis; the missing construction must determine its actual role. The research summary calls the flatness conclusion known while the shrinking/convergence mechanism remains unresolved.

Adjacent extracted material starts section 24 on scalar and mean curvature. It reports a doubling lemma [52]: a Riemannian \(Y\) with \(\operatorname{Sc}_Y\ge0\) and strictly mean-convex boundary has a continuous doubled metric on \(Y\cup_{\partial Y}Y\) approximable in \(C^0\) by \(C^2\) metrics of strictly positive scalar curvature, with \(C^2\) convergence away from the seam. The stated method smooths near the seam and then redistributes positivity conformally; compactness and mean-curvature conventions are not supplied. The following thin-space definition begins an infimum of boundary volumes over arbitrarily small neighborhoods and then ends in an ellipsis. This adjacent lemma and unfinished definition are not added as hypotheses of the shrinking conjecture.
\subsection{Short English statement}
In the missing regularized descent construction, prove diameter collapse, eventual smoothness of minimizing hypersurfaces, and bilipschitz approximation by flat tori with distortion tending to one. The source distinguishes this mechanism from its reported overtorical flatness consequence.
\subsection{Sources}
\begin{itemize}
\item Local UnsolvedMath registry, permanent ID \(6700065\), exact source code \texttt{catalogue entry}; printed scalar-curvature tag [?69], occurrence 68, ``Shrinking of Singularities.'' Proposer Misha Gromov, 2017.
\item Gromov, \emph{101 Questions, Problems and Conjectures around Scalar Curvature} (2017), recorded PDF page 64; \url{https://www.ihes.fr/~gromov/wp-content/uploads/2018/08/101-problemsOct1-2017.pdf}. Recorded access 2026-07-29; pdf-text-normalized extraction.
\item Source reference [49] for a \(C^0\) semicontinuity argument, [86,115] for possible descent/nonpositivity arguments, and [52] for the adjacent doubling lemma; bibliography entries are absent.
\item Supplied research summary names Schoen--Yau torus rigidity and relates the mechanism to Lohkamp/Schoen--Yau descent, without complete citations or precise hypotheses.
\end{itemize}
\end{document}
